\subsection{Purpose}
BBP is a software designed to simplify the life of everyone who goes on bike trips and wants to enjoy them without worrying about finding the best routes.\\
It is intended as a community for cyclists, in which every person can share information about their personal routes, so that everyone may take advantage of the information and experience of other bikers.
The application is intended for smartphone, where it can make use of GPS and sensor data to seamlessly reconstruct the path a cyclist has taken.

\subsection{Scope}

The scope of this Design Document (DD) is to define the technical architecture and the detailed design of the Best Bike Paths (BBP) system. 

It translates the requirements specified in the RASD into a modular, microservice-based structure, detailing the components, their interfaces, and the user interface that represents the implementation of the system. 

It represents the structure to respect in development to ensure that the system is scalable, maintainable, and that it correctly provides a reliable experience for users.

\subsection{Definitions, Acronyms, Abbreviations}

\begin{table}[H]
   \begin{tabular}{|p{3cm}|p{12cm}|}
   \hline
   \textbf{Acronym} & \textbf{Definition} \\ \hline
   BBP & Best Bike Paths application \\ \hline
   GPS & Global Positioning System \\ \hline
   UI & User Interface \\ \hline
   SP & Shared Phenomenon \\ \hline
   WP & World Phenomenon \\ \hline
   UC & Use Case \\ \hline
   DA & Domain Assumption \\ \hline
   G & Goal \\ \hline
   \end{tabular}
   \caption{Definitions, Acronyms and Abbreviations}
\end{table}

\subsection{Reference Documents}

\begin{itemize}
   \item Requirements analysis and specification document (RASD) of Best Bike Paths application
   \item Lecture slides and notes from the course Software Engineering 2 2025/2026
\end{itemize}

\subsection{Document Structure}

This document is organized into five main sections, each covering a specific aspect of the BBP (Best Bike Paths) system design:

\begin{itemize}
    \item \textbf{Section 1 - Introduction}: This section provides an overview of the document, defining its purpose and scope. 
    It also includes a list of definitions, acronyms, and abbreviations used throughout the text, along with references to related documents.
    \item \textbf{Section 2 - Architectural Design}: This section describes the system's high-level architecture. It details the microservice-based approach through a Component View and a Deployment View. 
    Furthermore, it includes a Runtime View where sequence diagrams illustrate how various use cases are executed across the components, and a dedicated subsection for Component Interfaces.
    \item \textbf{Section 3 - User Interface Design}: This section presents the design of the mobile application's interface. 
    It explains the layout and navigation flow, highlighting how the UI has been structured to satisfy the specific needs and goals of the end-users.
    \item \textbf{Section 4 - Requirements Traceability}: This section provides a mapping between the architectural components defined in this document and the functional requirements previously identified in the RASD (Requirements Analysis and Specification Document), 
    ensuring every requirement is addressed by the design.
    \item \textbf{Section 5 - Implementation, Integration, and Test Plan}: The final section outlines the development strategy. It identifies the order in which subcomponents will be implemented and describes the plan for integrating and testing them to ensure system reliability.
    \item \textbf{Section 6 - Effort Spent}: This section provides an overview of the effort spent by each team member during the creation of this document.
\end{itemize}